% Chapter 11: Technology Stack and Implementation
\chapter{Technology Stack and Implementation}
\label{chap:technology}

This chapter documents the concrete technologies used to implement the web spinner mechanics and how the theoretical constructs from Part II are realized in a working system.

\section{Runtime and Languages}
The reference implementation is written in JavaScript and runs on Node.js 22.18 or later:
\begin{itemize}
	\item Node.js runtime and JavaScript for server-side logic
	\item Standard Web APIs (URL, URLSearchParams, Fetch, Crypto) alongside Node.js APIs and npm packages
	\item npm scripts and Node.js utilities for development, testing, asset builds, and releases
	\item Containerization via a Dockerfile based on the official Node.js image
\end{itemize}

\subsection{Why Node.js for the \webspinner?}
Node.js provides a mature server runtime and package ecosystem while retaining the modern JavaScript and Web API features used by the \webspinner:
\begin{itemize}
	\item \textbf{JavaScript end-to-end:} one language for the spinner, utilities, and any shared logic with client code.
	\item \textbf{Web APIs in the server:} standard \texttt{URL}, \texttt{Request/Response}, \texttt{Fetch}, and \texttt{Crypto} reduce bespoke abstractions.
	\item \textbf{Package ecosystem:} npm dependencies provide database drivers, XML parsing, and WebSocket support.
	\item \textbf{Integrated tooling:} npm scripts and the built-in Node.js test runner support development and releases.
	\item \textbf{Simple deploys:} reproducible Docker images and the widely available Node.js runtime keep operations light.
	\item \textbf{Native ESM:} native modules and top-level \texttt{await} fit the compilation model used by \wbll and \wbpl tooling.
\end{itemize}

Other implementations are welcomed, provided they adhere to the specifications and interoperability guidelines set forth by the \organization{}.

\section{High-Level Architecture}
The spinner is a server that understands WBOL. On each request it:\footnote{See also the ``Paradigm'' section in the implementation README}
\begin{enumerate}
	\item Establishes or resumes a session (user, roles, locale, placeholders)
	\item Ensures the requested WBOL/Webbase is loaded into an in-memory tree
	\item Decides whether to respond with a resource directly or with a list of REST calls the client should make (async via WebSockets)
	\item Renders contents on demand using WBLL-driven layouts and returns HTML fragments/resources
\end{enumerate}

\section{Core Elements and Portal Tree}
The in-memory model mirrors the WBOL structure:
\begin{itemize}
	\item \textbf{STWPortal} (singleton root), \textbf{STWArea}, \textbf{STWPage}, \textbf{STWContent}
	\item All derive from an abstract \textbf{STWElement} providing identity, naming, localization, hierarchy, and export to WBOL
	\item A content type (e.g., Text, Table, Menus, Breadcrumbs, Calendar, Code Editor, ...), implemented under \texttt{stwContents/}, encapsulates data access and rendering concerns
\end{itemize}

\section{Rendering Pipeline: WBLL and WBPL}
Presentation is described with WBLL (Webbase Layout Language), interpreted by a layout interpreter:
\begin{itemize}
	\item WBLL strings are tokenized and validated; tokens drive generation of a specialized render function
	\item Token handlers cover inputs, lists, links, media, buttons/actions, and structural fragments; they build HTML using placeholders and field values
	\item WBPL expressions provide string interpolation and conditional/functional logic within layouts and settings
	\item Placeholders (e.g., \verb|@@name|) merge session, request, and record values
\end{itemize}

\section{Request Flow in Practice}
For a content render request:
\begin{enumerate}
	\item The session determines visibility (role-based) and language
	\item The content locates its WBLL layout for the current language
	\item Records are fetched (via a datasource or parameters); the first row and fields hydrate placeholders
	\item The compiled WBLL render function executes, producing the HTML body; optional header/footer wrappers apply
\end{enumerate}

\section{Security, Localization, and State}
\begin{itemize}
	\item \textbf{Visibility}: role-based flags inherited along the element tree control exposure of nodes
	\item \textbf{Localization}: localized properties (names, slugs, messages) are resolved through the session
	\item \textbf{State}: per-session placeholders and content-level settings influence rendering and actions
\end{itemize}

\section{Build, Tooling, and Deployment}
\begin{itemize}
	\item npm scripts run Node.js tasks to merge and minify static assets (e.g., CSS and JavaScript) for the \texttt{public/} client
	\item Tests: parsing and evaluation tests for WBPL ensure correctness of expressions and escaping
	\item Docker: a Node.js image installs npm dependencies and ships a non-root runtime image
\end{itemize}

\section{Where the Mechanics Live (Guide to Source)}
The following folders in the reference implementation contain the mechanics described above:
\begin{itemize}
	\item \texttt{stwElements/}: \textit{STWElement}, \textit{STWPortal}, \textit{STWArea}, \textit{STWPage}, \textit{STWContent}
	\item \texttt{stwContents/}: concrete contents and \textit{WBLL} interpreter (layout parsing/rendering)
	\item \texttt{tests/}: WBPL and layout-related unit tests
	\item \texttt{public/}: client-side scripts, styles, and SPA shell
	\item \texttt{tasks/}: Node.js development and build utilities (e.g., minification, CSS merge)
\end{itemize}

\section{Example: From WBOL to HTML}
At a glance:
\begin{enumerate}
	\item WBOL defines the portal tree; the spinner builds an in-memory model at startup/load
	\item A user navigates to a page; the spinner locates the route and the associated contents
	\item Each content loads data (if needed), prepares placeholders, and renders its WBLL layout
	\item The server responds with an HTML fragment or instructs the client to fetch multiple fragments via REST/WebSockets
\end{enumerate}

\noindent This chapter has bridged theory and implementation, showing how the abstract mechanics of the Web Spinner are realized in a modern technology stack. With the platform's construction now fully detailed—from its languages to its runtime and development tools—our focus shifts from building the framework to using it.

The next part of this book begins the practical guide for developers, covering the development models and methodologies for designing and building effective portals with the platform we have just described.
